\documentclass[11pt,a4paper]{report}
\usepackage[margin=2.4cm]{geometry}\usepackage{booktabs,tabularx}\usepackage[table]{xcolor}\usepackage{enumitem,fancyhdr}
\definecolor{blue}{HTML}{244A73}\pagestyle{fancy}\fancyhf{}\lhead{Orchard Analytics | Project report}\rhead{\thepage}\renewcommand{\headrulewidth}{.4pt}
\begin{document}\begin{titlepage}\centering\vspace*{2cm}{\color{blue}\Huge\bfseries ORCHARD ANALYTICS\par}\vspace{.5cm}{\Large Internship and company project report\par}\vspace{2cm}{\huge Customer-support forecasting pilot\par}\vfill\begin{tabular}{rl}Author:& Maya Fernandes\\Placement:& Data Operations, Spring 2025\\Supervisor:& Eli Morgan\\Date:& 30 June 2025\end{tabular}\vfill{\small A complete starter with scope, evidence, delivery log and reflection.}\end{titlepage}
\tableofcontents\chapter{Executive summary} This project tested whether a lightweight forecast could help the support team schedule Monday staffing. A cleaned six-month ticket history, a transparent seasonal baseline and a short handover guide were delivered. The pilot reduced the largest weekly staffing miss from 18 tickets to 7 in the final test month; the result is encouraging but not sufficient to automate rota decisions.
\chapter{Context and scope}\section{Business question} How can the team anticipate the next week’s inbound ticket volume without adding a fragile production dependency?\section{Success criteria}\begin{itemize}[leftmargin=*]\item produce a forecast before the Monday planning meeting;\item show the assumptions behind every number;\item leave a repeatable handover, including known limitations;\item avoid exposing customer-identifying fields in the analysis.\end{itemize}
\chapter{Delivery record}\begin{tabularx}{\textwidth}{p{.18\textwidth}p{.3\textwidth}X}\toprule Period & Activity & Evidence / outcome\\\midrule Weeks 1--2 & Data audit & 18,420 tickets; duplicate and reopened-ticket rules documented.\\Weeks 3--4 & Baseline model & Weekly seasonal mean, compared with a naive last-week baseline.\\Weeks 5--6 & Pilot & Forecast shared with two planners; feedback captured after each meeting.\\Weeks 7--8 & Handover & One-page runbook, validation checklist and risk log delivered.\\\bottomrule\end{tabularx}
\chapter{Findings}\section{What changed} The baseline was most useful as a conversation starter: planners could see the effect of a public-holiday adjustment and challenge it before committing overtime. It was less useful during a product outage, which created a spike outside historical patterns.\section{Risks and limitations}\begin{enumerate}[leftmargin=*]\item The dataset represents one product and should not be extrapolated to the whole company.\item Forecast accuracy was measured on four weeks only.\item The pipeline depends on a manually exported CSV and needs an owner before reuse.\end{enumerate}
\chapter{Reflection and next steps}\textbf{Keep:} a small model, visible assumptions and a review conversation. \textbf{Improve:} add an outage indicator, automate the export and pre-register the evaluation window. My main learning was that a usable business deliverable includes adoption instructions and ownership, not only an analysis.
\appendix\chapter{Handover checklist}\begin{itemize}[leftmargin=*]\item Export the prior week’s ticket count.\item Check missing dates and duplicate IDs.\item Record any outage, campaign or policy change.\item Compare forecast and actual after the week closes.\item Update the assumptions and notify the supervisor of material drift.\end{itemize}\end{document}
